% =============================================================================
% Day 1 Course Overview — ML Systems Lecture Slides
% =============================================================================
% This deck is presented on the FIRST DAY of class to introduce students
% to the course identity, philosophy, analytical frameworks, and logistics.
%
% IMAGES NEEDED (generate SVGs, convert to PDF):
%   - course-roadmap.pdf       - deployment-spectrum.pdf
%   - frameworks-overview.pdf  - lighthouse-models.pdf
%   - five-percent-overview.pdf
% =============================================================================
\documentclass[aspectratio=169, 12pt]{beamer}
\usepackage{../../assets/beamerthememlsys}

\mlsyssetup{
  volume       = {Volume I},
  chapter      = {Course Overview},
  logo         = {../../assets/img/logo-mlsysbook.png},
  instlogo     = {../../assets/img/logo-harvard.png},
  chaptertitle = {Introduction to ML Systems},
}

% --- Fonts ---
\usepackage[T1]{fontenc}
\usepackage[scaled=0.9]{helvet}
\usepackage{courier}
\renewcommand{\familydefault}{\sfdefault}

% --- Packages ---
\usepackage{booktabs}
\usepackage{amsmath}

% --- Image paths ---
\graphicspath{
  {images/}
}

% --- Chapter-specific macros ---
\newcommand{\DAM}{D\raisebox{0.08em}{\tiny$\bullet$}A\raisebox{0.08em}{\tiny$\bullet$}M}

% --- Helper: safe image include ---
\newcommand{\safeimg}[2][width=\textwidth,keepaspectratio]{%
  \IfFileExists{images/#2}{\includegraphics[#1]{#2}}{%
    \IfFileExists{#2}{\includegraphics[#1]{#2}}{%
      \fbox{\parbox[c][2.5cm][c]{0.85\linewidth}{\centering\footnotesize\textcolor{midgray}{[Missing image]}}}%
    }%
  }%
}

% --- Section count for navigation (must match actual \section{} count) ---
\setcounter{mlsystotalsections}{7}

\title{Course Overview}
\author{Vijay Janapa Reddi}
\institute{Harvard University}
\date{}

\begin{document}

% =============================================================================
% TITLE SLIDE
% =============================================================================
\mlsystitle{Course Overview}{The Physics of AI Engineering}{cover_course_overview.png}

% =============================================================================
% VISUAL LANGUAGE & ORIENTATION
% =============================================================================

\begin{frame}{Visual Language}
\note{
% [1 min]
% -- LINK: This is the first slide students see. Before any content, establish
% the visual grammar so color carries meaning throughout the semester.
%
% -- NARRATE: Walk through each card. ``Blue means compute---anytime you see
% blue in a diagram, think GPU ops, matrix multiplies, forward pass. Green is
% data flow---memory, caches, healthy paths. Orange is routing---schedulers,
% load balancers. Red is cost, error, or bottleneck.'' Point to each card as
% you name it.
%
% -- ENGAGE: Before revealing the colors, ask: ``What color would you use for
% memory bandwidth?'' Expected: Green (data flow). Students who say blue will
% need the distinction: blue is arithmetic, green is movement.
%
% -- WARN: Students often ignore color coding and lose the diagram's diagnostic
% signal all semester. Emphasize: the color is not decoration, it is the
% bottleneck indicator. If something is red, that is the constraint.
%
% -- FLEX: [CORE] Show on Day 1 and again briefly on Day 2.
% IF SHORT: Display the slide but do not narrate---students can read the cards.
}

\small
Throughout this course, colors carry meaning:

\vspace{0.3cm}
\begin{columns}[T]
  \begin{column}{0.45\textwidth}
    \begin{mlsyscard}{computestroke}
    \textbf{Blue} --- Compute / Processing\\
    {\footnotesize GPU ops, forward/backward pass, inference}
    \end{mlsyscard}
    \vspace{0.15cm}
    \begin{mlsyscard}{datastroke}
    \textbf{Green} --- Data / Memory\\
    {\footnotesize Data flow, caches, healthy paths}
    \end{mlsyscard}
  \end{column}
  \begin{column}{0.45\textwidth}
    \begin{mlsyscard}{routingstroke}
    \textbf{Orange} --- Routing / Scheduling\\
    {\footnotesize Load balancers, batch windows}
    \end{mlsyscard}
    \vspace{0.15cm}
    \begin{mlsyscard}{errorstroke}
    \textbf{Red} --- Error / Cost / Bottleneck\\
    {\footnotesize Loss, decode phase, waste}
    \end{mlsyscard}
  \end{column}
\end{columns}
\end{frame}

\begin{frame}{How to Read This Textbook}
\note{
% [2 min]
% -- LINK: Students just learned the visual language. This slide orients them
% to the textbook structure they will use all semester.
%
% -- NARRATE: ``The textbook has two volumes. Volume I---this course---covers
% single-machine systems: 1--8 GPUs, shared memory, one node. Volume II covers
% multi-machine distributed systems---that is a follow-on course. Volume I has
% four parts: Foundations of ML Systems, Development and Training,
% Optimization and Acceleration, and Deployment and Operations.
% Five lighthouse models recur in every chapter as running examples.''
%
% -- ENGAGE: ``How many of you have used a two-volume textbook before?'' Show
% of hands.
%
% -- WARN: Students may skip Part I thinking they ``already know ML.'' Correct
% framing: Part I establishes the analytical vocabulary that every later chapter
% depends on.
%
% -- FLEX: [CORE] Orientation slide for Day 1.
% IF SHORT: Show for 60 seconds, highlight the four parts only.
}

\footnotesize
\begin{columns}[T]
  \begin{column}{0.48\textwidth}
    \textbf{Two-volume structure:}

    \vspace{0.1cm}
    \textcolor{computestroke}{\textbf{Volume I}} (this course):\\
    {\scriptsize Single-machine systems, 1--8 GPUs, shared memory}

    \vspace{0.1cm}
    \textcolor{datastroke}{\textbf{Volume II}} (follow-on):\\
    {\scriptsize Distributed systems, multi-node, fleet-scale}

    \vspace{0.15cm}
    \textbf{Four Parts in Volume I:}
    \begin{enumerate}\setlength\itemsep{0pt}
      \item {\scriptsize \textbf{Foundations of ML Systems} --- frameworks \& physics}
      \item {\scriptsize \textbf{Development and Training} --- data, training, architectures}
      \item {\scriptsize \textbf{Optimization and Acceleration} --- compression \& hardware}
      \item {\scriptsize \textbf{Deployment and Operations} --- serving, ops, ethics}
    \end{enumerate}
  \end{column}
  \begin{column}{0.48\textwidth}
    \vspace{0.1cm}
    \begin{mlsyscard}{crimson}
    \textbf{Lighthouse Models}\\[0.1cm]
    {\scriptsize Five models recur in every chapter:}
    \begin{itemize}\setlength\itemsep{0pt}
      \item {\scriptsize ResNet-50 (compute-bound)}
      \item {\scriptsize GPT-2/Llama (memory-BW-bound)}
      \item {\scriptsize MobileNetV2 (latency-bound)}
      \item {\scriptsize DLRM (memory capacity-bound)}
      \item {\scriptsize KWS (power-bound)}
    \end{itemize}
    {\scriptsize Same concept, different constraint regime.}
    \end{mlsyscard}
  \end{column}
\end{columns}

\end{frame}

% =============================================================================
% WELCOME & COURSE IDENTITY
% =============================================================================

\begin{frame}{Welcome}
\note{
% [2 min]
% -- LINK: Students are arriving on Day 1 with varying backgrounds. Start with
% what they all share: most have trained a model but few have deployed one.
%
% -- NARRATE: ``Welcome. Raise your hand if you have trained a model.'' (most
% hands go up) ``Now keep your hand up if you have deployed one to production.''
% (most hands drop) ``That gap---training vs.\ shipping---is this entire course.
% This is not an ML algorithms class. This is about the systems that make ML
% work: memory, bandwidth, power, latency, and the physics behind every design
% decision.''
%
% -- ENGAGE: Ask for a show of hands. Wait for the drop---that visual moment
% is the emotional hook. Expected response: most hands stay up for training,
% most fall for deployment. Follow up: ``The gap between those two moments
% is why this course exists.''
%
% -- WARN: Students may assume this is another deep-learning or PyTorch course.
% Correct framing immediately: the bottleneck is systems, not algorithms.
%
% -- FLEX: [CORE] Sets the emotional contract for the semester.
% IF SHORT: Cut the hand-raise; just state the training-deployment gap directly.
}

\centering
\vspace{0.8cm}
{\Large\bfseries Introduction to Machine Learning Systems}

\vspace{0.4cm}
{\large Volume I: Foundations for Single-Machine ML Systems}

\vspace{0.6cm}
{\normalsize\textcolor{crimson}{\textbf{``The Physics of AI Engineering''}}}

\vspace{0.6cm}
{\small AI is not magic --- it is infrastructure, and infrastructure has laws.}

\end{frame}

% =============================================================================
\section{Why This Course}
% =============================================================================

\begin{frame}{The Gap Between ML Research and Production}
\note{
% [3 min]
% -- LINK: Students just heard ``this is a systems course, not an algorithms
% course.'' This slide gives the quantitative evidence for WHY.
%
% -- NARRATE: ``60 to 85 percent of ML projects never reach production. Let
% that sink in. Look at this table---research uses a static dataset, a single
% GPU, and optimizes one metric. Production uses a shifting data stream, a
% fleet of heterogeneous hardware, and must hit accuracy AND latency AND cost
% targets simultaneously. That gap is not fixed by a better optimizer.''
%
% -- ENGAGE: ``Why do you think most ML projects fail? Write one reason.''
% Give 30 seconds. Cold-call 2 students.
% Expected: ``data quality,'' ``hardware limits.'' Surprise answer: systems.
%
% -- WARN: Students assume failures are algorithmic (``bad model''). Correct
% framing: the bottleneck is infrastructure---data pipelines, serving, monitoring.
% IF STUCK: Point to the ``90\% of time goes to data + infrastructure'' bullet.
%
% -- FLEX: [CORE] This slide motivates the entire semester.
% IF AHEAD: ``What percentage of engineering time goes to the model itself?''
}

\footnotesize
\begin{columns}[T]
  \begin{column}{0.55\textwidth}
    A Kaggle notebook is not a production system:

    \vspace{0.15cm}
    \begin{itemize}\setlength\itemsep{0pt}
      \item \textbf{60--85\%} of ML projects fail to reach production
      \item 90\% of engineering time goes to data + infrastructure
      \item Model code is $\approx$\textbf{5\%} of a production system
    \end{itemize}

    \vspace{0.15cm}
    \begin{mlsyscard}{errorstroke}
    The bottleneck is \alert{systems}, not algorithms.
    \end{mlsyscard}
  \end{column}
  \begin{column}{0.42\textwidth}
    \vspace{0.15cm}
    \renewcommand{\arraystretch}{1.15}
    \begin{tabular}{@{}l@{\hspace{4pt}}l@{\hspace{4pt}}l@{}}
      \toprule
      & \textbf{Research} & \textbf{Production} \\
      \midrule
      Data       & Static set   & Shifting stream \\
      Hardware   & Given GPU    & Fleet of tiers \\
      Success    & Accuracy     & Accuracy + SLA \\
      Lifespan   & One paper    & Years \\
      \bottomrule
    \end{tabular}
  \end{column}
\end{columns}
\mlsyscite{Paleyes et al., ACM Computing Surveys 2022}

\end{frame}

\begin{frame}{The 5\% Problem}
\note{
% [2 min]
% -- LINK: The previous slide said ``the bottleneck is systems.'' This diagram
% shows exactly what those systems look like.
%
% -- NARRATE: Point to the tiny center box: ``That is the ML model code---about
% 5 percent. Everything around it---data pipelines, feature stores, serving
% infrastructure, monitoring, configuration---is what this course teaches.
% Sculley et al.\ called this `hidden technical debt.' ''
%
% -- ENGAGE: ``Look at the diagram. What is the biggest box?'' Give 15 seconds.
% Expected answer: data collection or configuration. Both are valid---the point
% is that neither is the model.
%
% -- WARN: Students equate ``ML'' with ``the model.'' Correct framing: the
% model is 5\%; the other 95\% is systems engineering that determines whether
% the model ever reaches a user.
%
% -- FLEX: [CORE] Foundational mental model for the course.
% IF SHORT: Skip the question; just narrate the diagram for 90 seconds.
}

% --- Layout: FULL-WIDTH IMAGE ---
\centering
\safeimg[width=0.88\textwidth,height=4.8cm,keepaspectratio]{five-percent-overview.pdf}

\vspace{0.1cm}
{\small This course teaches the \textbf{95\%} --- the systems engineering that makes ML work.}
\mlsyscite{Sculley et al., NeurIPS 2015}

\end{frame}

% =============================================================================
% WHY ML SYSTEMS ARE DIFFERENT --- THREE THROUGH-LINE SLIDES
% =============================================================================

\begin{frame}{Why ML Systems Are Different}
\note{
% [2 min]
% -- LINK: The 5\% diagram showed that the bottleneck is systems. This slide
% explains *why* systems engineering for ML is genuinely different from
% traditional software engineering.
%
% -- NARRATE: ``When a traditional program has a bug, it crashes. You get a
% stack trace in milliseconds. When an ML system has a bug, it keeps running,
% returns predictions, and looks healthy---while silently becoming wrong.
% A sentiment model trained on last year's data drops from 94\% to 78--82\%
% accuracy in production as users adopt new slang and context. No exception
% thrown. No alert. The code executes correctly. The problem is in the data.''
%
% -- ENGAGE: ``What would you do if your ML system's accuracy dropped 5 points
% with no code changes?'' Give 30 seconds. Expected: check the data, check the
% distribution, look at monitoring dashboards. Follow-up: ``What if you had no
% monitoring?'' Expected: you would not know.
%
% -- WARN: Students assume ``no bugs = working correctly.'' This is the most
% dangerous misconception in ML engineering. Correct framing: ML systems degrade
% through data drift even when every line of code is correct. This is why the
% degradation equation and continuous monitoring are engineering requirements,
% not optional nice-to-haves.
%
% -- FLEX: [CORE] The foundational argument for why this course exists.
% IF SHORT: State the silent/loud contrast in one sentence and move on.
}

\footnotesize
\begin{columns}[T]
  \begin{column}{0.50\textwidth}
    \textbf{Traditional software:}

    \vspace{0.08cm}
    \begin{itemize}\setlength\itemsep{0pt}
      \item Code bug $\to$ crash (loud failure)
      \item Stack trace in milliseconds
      \item Debugging: trace the code
      \item Fixed code stays fixed
    \end{itemize}

    \vspace{0.1cm}
    \begin{mlsyscard}{computestroke}
    {\footnotesize A function that computed correctly yesterday
    computes correctly today.}
    \end{mlsyscard}
  \end{column}
  \begin{column}{0.47\textwidth}
    \textbf{ML systems:}

    \vspace{0.08cm}
    \begin{itemize}\setlength\itemsep{0pt}
      \item Data bug $\to$ wrong prediction (silent failure)
      \item No exception, no alert
      \item Debugging: trace the \emph{data}
      \item Model silently decays as world changes
    \end{itemize}

    \vspace{0.05cm}
    \begin{mlsyscard}{errorstroke}
    {\footnotesize Sentiment model: 94\% on test set $\to$
    78--82\% in production as users adopt new slang.
    \textbf{Code unchanged.}}
    \end{mlsyscard}
  \end{column}
\end{columns}

\vspace{0.05cm}
\mlsysconcept{Implication:}{You cannot test your way to reliability. You must \emph{monitor} it.}

\end{frame}

\begin{frame}{Software 2.0 --- Weights Are the New Source Code}
\note{
% [2 min]
% -- LINK: The previous slide showed that ML systems fail through data bugs,
% not code bugs. This slide explains *why*---the program logic lives in
% the weights, not in the code.
%
% -- NARRATE: ``Andrej Karpathy named this shift in 2017: Software 1.0 is
% explicit instructions written by a programmer. Software 2.0 is a neural
% network whose behavior is defined by weights learned from data.
% The code that trains a model is only a few thousand lines.
% The actual program is 1.8 trillion floating-point numbers in a GPT-4-class
% model. You cannot read it. You cannot patch it. You can only retrain it.''
% ANALOGY: ``Traditional debugging is like editing source code with a text
% editor. ML debugging is like finding the one bad ingredient in a batch of
% 1.8 trillion cookies---by baking a new batch.''
%
% -- ENGAGE: ``Where does the `program logic' live in a GPT-4-class model?''
% Expected: in the weights (not in the Python code). Follow-up:
% ``What does version control mean when your `source code' is a 3.6 TB file?''
%
% -- WARN: Students think of models as black boxes. Correct framing:
% the weights ARE the program---and the entire ML systems engineering
% discipline exists because those weights are produced by, and sensitive to,
% the data pipeline, the training infrastructure, and the serving environment.
%
% -- FLEX: [CORE] The paradigm-shift framing that motivates the 95\% problem.
% IF SHORT: State the weights-as-source-code analogy and move to the next slide.
}

\small
\begin{columns}[T]
  \begin{column}{0.48\textwidth}
    \begin{mlsyscard}{computestroke}
    \textbf{Software 1.0}\\[0.1cm]
    {\footnotesize
    \textbf{Program}: explicit instructions\\
    \textbf{Bug}: trace to a line of code\\
    \textbf{Update}: edit the source file\\
    \textbf{Test}: deterministic paths\\
    }
    \end{mlsyscard}
  \end{column}
  \begin{column}{0.48\textwidth}
    \begin{mlsyscard}{errorstroke}
    \textbf{Software 2.0}\\[0.1cm]
    {\footnotesize
    \textbf{Program}: learned weights\\
    \textbf{Bug}: trace to the training data\\
    \textbf{Update}: retrain on new data\\
    \textbf{Test}: monitor production distributions\\
    }
    \end{mlsyscard}
  \end{column}
\end{columns}

\vspace{0.2cm}
\mlsysalert{Systems consequence:}{A GPT-4-class model has $\sim$1.8T parameters. The ``program'' is 3.6 TB of weights. Debugging means auditing the data pipeline, not the Python code.}
\mlsyscite{Karpathy, 2017}

\end{frame}

\begin{frame}{The Bitter Lesson: Scale Beats Expertise}
\note{
% [2 min]
% -- LINK: Software 2.0 told us that behavior comes from data and scale,
% not from hand-coded logic. The Bitter Lesson explains *why* this pattern
% holds across all of AI history---and why it forces ML systems engineering
% to exist.
%
% -- NARRATE: ``Richard Sutton looked at 70 years of AI research and found
% one consistent pattern: every time a paradigm relied on encoding human
% expertise, it eventually lost to a simpler method that could absorb more
% compute. Rule-based AI lost to statistical ML. Hand-crafted features lost
% to learned representations. The lesson is bitter because it contradicts
% our intuition that smarter hand-engineering should win. It never does,
% at sufficient scale. The implication for systems engineers: the discipline
% that moves compute efficiently determines AI progress.''
%
% -- ENGAGE: ``If scale always wins, why does systems engineering matter?''
% Expected: because scaling requires infrastructure---data pipelines, efficient
% training, hardware optimization. Without systems engineering, scale is
% physically impossible regardless of the algorithm.
%
% -- WARN: Students may think the lesson implies algorithms do not matter.
% Correct framing: the lesson is about the *relative* importance of scale vs.
% hand-coded knowledge. Algorithmic insight still matters---it just matters
% less than compute-efficient infrastructure at the frontier.
%
% -- FLEX: [CORE] The historical and motivational argument for the whole course.
% IF SHORT: State the two-line version: ``Scale beats expertise. Systems
% engineering is what makes scale possible.''
}

\footnotesize
\begin{columns}[T]
  \begin{column}{0.55\textwidth}
    \textbf{Pattern across 70 years of AI:}

    \vspace{0.08cm}
    \begin{itemize}\setlength\itemsep{0pt}
      \item \textbf{1950s--70s:} Symbolic AI --- logic rules, expert knowledge
      \item \textbf{1980s--90s:} Feature engineering --- human-crafted representations
      \item \textbf{2000s--10s:} Statistical ML --- learned features, manual tuning
      \item \textbf{2012--now:} Deep learning + massive compute \textcolor{computestroke}{\textbf{win}}
    \end{itemize}

    \vspace{0.08cm}
    \mlsysalert{The Lesson:}{``General methods that leverage computation consistently outperform approaches that encode human expertise.'' --- Sutton, 2019}
  \end{column}
  \begin{column}{0.42\textwidth}
    \vspace{0.08cm}
    \begin{mlsyscard}{crimson}
    {\footnotesize \textbf{Systems implication:} If scale determines AI progress, then the systems that move compute efficiently---data pipelines, training infrastructure, hardware optimization---are the limiting factor.

    \vspace{0.08cm}
    \textbf{This course teaches how to build those systems.}}
    \end{mlsyscard}
  \end{column}
\end{columns}
\mlsyscite{Sutton, ``The Bitter Lesson,'' 2019}

\end{frame}

% =============================================================================
\section{The Core Thesis}
% =============================================================================

\mlsysfocus{The Core Thesis}{%
Constraints drive architecture.\\[0.3cm]
You don't choose a Transformer because it's trendy ---\\
you choose it because of how it parallelizes on real silicon.%
}

\begin{frame}{AI Is Infrastructure}
\note{
% [3 min]
% -- LINK: The Core Thesis focus slide just said ``constraints drive
% architecture.'' This slide names the four specific physical constraints.
%
% -- NARRATE: Walk through bullets top to bottom. ``Memory bandwidth limits
% how fast data reaches the processor---this is why GPT-4 inference is slow
% even on powerful GPUs. Power budget limits where a model can run---a phone
% cannot sustain 700 watts. Speed of light limits latency---a self-driving car
% cannot wait 50 ms for a cloud round trip. Thermodynamics limits compute per
% rack---you cannot cool infinite GPUs in a data center.''
% ANALOGY: ``Physics is to ML systems what gravity is to bridges. You can
% build creative bridges, but none of them ignore gravity.''
%
% -- ENGAGE: ``Why can't we run GPT-4 on a smartphone?'' Cold-call one student.
% Expected: ``not enough memory.'' Deepen: ``How much memory does it need?
% 3.6 TB at FP16 (1.8T params $\times$ 2 bytes). A phone has 8 GB.
% That is a 450x gap---physics, not software.''
%
% -- WARN: Students think hardware limitations are temporary (``next year's chip
% will fix it''). Correct framing: these are physical laws, not engineering gaps.
% Memory bandwidth grows ~20\%/yr; compute demand grows 7x faster than Moore's Law.
%
% -- FLEX: [CORE] Philosophical anchor for the course.
% IF AHEAD: ``Which constraint is hardest to overcome with engineering?''
}

\small
\begin{columns}[T]
  \begin{column}{0.55\textwidth}
    \textbf{Physical laws govern ML systems:}

    \vspace{0.1cm}
    \begin{itemize}\setlength\itemsep{0pt}
      \item \textcolor{computestroke}{\textbf{Memory bandwidth}} limits how fast data reaches the processor
      \item \textcolor{routingstroke}{\textbf{Power budget}} limits where a model can run
      \item \textcolor{errorstroke}{\textbf{Speed of light}} limits latency over distance
      \item \textcolor{datastroke}{\textbf{Thermodynamics}} limits compute per rack
    \end{itemize}

    \vspace{0.05cm}
    \mlsysconcept{Implication:}{Physics, not elegance, determines what ships.}
  \end{column}
  \begin{column}{0.42\textwidth}
    \vspace{0.1cm}
    \begin{mlsyscard}{crimson}
    {\footnotesize \textbf{Example:} GPT-4 costs $\sim$\$100M to train --- moving 1.8T parameters through memory thousands of times costs that much in electricity and silicon time.}
    \end{mlsyscard}

    \vspace{0.1cm}
    \begin{mlsyscard}{datastroke}
    {\footnotesize \textbf{Example:} A self-driving car cannot call the cloud --- speed of light adds 50+ ms. Physics forces the model onto the edge.}
    \end{mlsyscard}
  \end{column}
\end{columns}

\end{frame}

% =============================================================================
\section{Three Frameworks}
% =============================================================================

\begin{frame}{Three Analytical Frameworks}
\note{
% [2 min]
% -- LINK: Students now know that physical constraints drive architecture.
% These three frameworks are the tools for reasoning about those constraints.
%
% -- NARRATE: Point to each framework in the diagram. ``D-A-M tells you WHERE
% the bottleneck is. The Iron Law tells you HOW LONG an operation takes. The
% Degradation Equation tells you WHEN a model will fail. These three tools
% recur in every single chapter.''
%
% -- ENGAGE: ``What question does each framework answer? Write one word per
% framework.'' 30 seconds. Cold-call one student.
% Expected: Where/How long/When---or close variants.
%
% -- WARN: Students will try to memorize the equations without understanding
% what question each answers. Correct framing: frameworks are diagnostic
% tools, not formulas to plug numbers into.
%
% -- FLEX: [CORE] Preview only---do not go deep. Chapter 1 covers each.
% IF SHORT: Just name the three frameworks and move on (60 seconds).
}

% --- Layout: FULL-WIDTH IMAGE ---
\centering
\safeimg[width=0.92\textwidth,height=4.8cm,keepaspectratio]{frameworks-overview.pdf}

\vspace{0.1cm}
{\small Three tools you will use in \textbf{every chapter} of this course.}

\end{frame}

\begin{frame}{The \DAM{} Taxonomy}
\note{
% [2 min]
% -- LINK: The previous diagram named D-A-M as one of three frameworks.
% Now we unpack it briefly.
%
% -- NARRATE: ``Every ML system sits at the intersection of three axes.
% Data: how much, how fast can we move it. Algorithm: how many operations,
% what parallelism pattern. Machine: what silicon, what memory hierarchy.
% The diagnostic question is always the same: which axis is the bottleneck?
% Optimizing one axis shifts pressure to the others---they are coupled.''
% Point to the table: ``Notice the units. These become the Iron Law variables.''
%
% -- ENGAGE: ``If you double the training dataset, which other axis feels
% the pressure?'' Expected: Machine (need more bandwidth or compute time).
%
% -- WARN: Students treat the three axes as independent knobs. Correct framing:
% they are interdependent---doubling data volume requires proportionally more
% bandwidth or longer training time.
%
% -- FLEX: [CORE] But keep it brief---Chapter 1 goes deep.
% IF SHORT: Show the slide for 60 seconds, skip the engage question.
}

\small
\begin{columns}[T]
  \begin{column}{0.55\textwidth}
    Every ML system has three interdependent axes:

    \vspace{0.2cm}
    \textcolor{datastroke}{\textbf{Data}} --- Volume, quality, distribution\\
    {\footnotesize How much data, how fast can we move it?}

    \vspace{0.1cm}
    \textcolor{computestroke}{\textbf{Algorithm}} --- Architecture, operation count\\
    {\footnotesize How many FLOPs, what parallelism pattern?}

    \vspace{0.1cm}
    \textcolor{routingstroke}{\textbf{Machine}} --- Hardware throughput, memory\\
    {\footnotesize What silicon, what memory hierarchy?}

    \vspace{0.15cm}
    \mlsysalert{Key:}{Optimizing one axis shifts pressure to another.}
  \end{column}
  \begin{column}{0.42\textwidth}
    \vspace{0.3cm}
    \renewcommand{\arraystretch}{1.2}
    {\footnotesize
    \begin{tabular}{@{}lll@{}}
      \toprule
      \textbf{Axis} & \textbf{Variable} & \textbf{Unit} \\
      \midrule
      Data      & $D_{\text{vol}}$, $BW$ & GB, GB/s \\
      Algorithm & $O$ & FLOPs \\
      Machine   & $R_{\text{peak}}$ & FLOP/s \\
      \bottomrule
    \end{tabular}
    }

    \vspace{0.2cm}
    \begin{mlsyscard}{crimson}
    {\footnotesize \textbf{Diagnostic:} ``Which axis is the bottleneck?'' answers most ML systems design questions.}
    \end{mlsyscard}
  \end{column}
\end{columns}

\end{frame}

\begin{frame}{The Iron Law of ML Systems}
\note{
% [2 min]
% -- LINK: D-A-M tells you where the bottleneck is. The Iron Law quantifies
% how long each axis takes in seconds.
%
% -- NARRATE: Point to each term. ``Data term: bytes divided by bandwidth
% gives seconds. Compute term: FLOPs divided by peak rate times efficiency
% gives seconds. Overhead: orchestration tax, also seconds. You add three
% times and the slowest one dominates. This is dimensional analysis---if
% your units do not resolve to seconds, the equation is wrong.''
%
% -- ENGAGE: ``For a phone camera app classifying a photo, which term
% dominates?'' Give 20 seconds. Expected: Data term (reading the image from
% memory) or Overhead (framework launch cost). Accept either with reasoning.
%
% -- WARN: Students try to add FLOPs to bytes. Correct framing: every term
% must resolve to seconds before you can compare or add them.
%
% -- FLEX: [CORE] Preview only---worked examples come in Chapter 1.
% IF SHORT: State the equation, emphasize ``slowest term dominates,'' move on.
}

\small
$$T_{\text{total}} = \underbrace{\dfrac{D_{\text{vol}}}{BW}}_{\text{Data}} +
\underbrace{\dfrac{O}{R_{\text{peak}} \cdot \eta}}_{\text{Compute}} +
\underbrace{L_{\text{lat}}}_{\text{Overhead}}$$

\vspace{0.1cm}
\begin{columns}[T]
  \begin{column}{0.30\textwidth}
    \centering
    \textcolor{datastroke}{\textbf{Data}} $\frac{D_{\text{vol}}}{BW}$\\[0.05cm]
    {\scriptsize Bytes / (Bytes/s) = \textbf{s}}
  \end{column}
  \begin{column}{0.30\textwidth}
    \centering
    \textcolor{computestroke}{\textbf{Compute}} $\frac{O}{R_p \cdot \eta}$\\[0.05cm]
    {\scriptsize FLOPs / (FLOPs/s) = \textbf{s}}
  \end{column}
  \begin{column}{0.30\textwidth}
    \centering
    \textcolor{errorstroke}{\textbf{Overhead}} $L_{\text{lat}}$\\[0.05cm]
    {\scriptsize Orchestration tax = \textbf{s}}
  \end{column}
\end{columns}

\vspace{0.15cm}
\centering
\mlsysconcept{Key insight:}{Every term resolves to \textbf{seconds}. The \textbf{slowest term dominates}.}

\end{frame}

\begin{frame}{The Degradation Equation}
\note{
% [2 min]
% -- LINK: The Iron Law measures performance at a point in time. The
% Degradation Equation measures how performance decays over time.
%
% -- NARRATE: ``Accuracy at time t approximately equals initial accuracy minus
% lambda times the distribution divergence. Lambda is the model's sensitivity
% to distribution shift---an architecture-dependent constant. D is the
% divergence measure (KL, total variation, or Wasserstein) quantifying how
% far live data has drifted from training data. Look at the example: a
% recommendation system starts at 85\% and drops to 79.2\% in 6 months.
% No code changed. No bugs. The world changed.
% The engineering response: set a retraining trigger at a threshold.''
%
% -- ENGAGE: ``How would you know your model is getting worse if no code
% changed and no one filed a bug?'' Give 20 seconds. Cold-call one student.
% Expected: monitoring accuracy metrics over time. Deepen: ``What if you
% do not have ground truth labels in real time?''
%
% -- WARN: Students assume ``no bugs = working correctly.'' Correct framing:
% ML systems degrade through data drift even when code is untouched.
%
% -- FLEX: [CORE] Third framework preview.
% IF SHORT: Show equation, state the rec-system example, move on.
}

\small
$$\text{Accuracy}(t) \approx \text{Accuracy}_0 - \lambda \cdot \mathcal{D}(P_t \lVert P_0)$$

\vspace{0.15cm}
\begin{columns}[T]
  \begin{column}{0.50\textwidth}
    {\footnotesize
    \begin{tabular}{@{}ll@{}}
      $\text{Accuracy}(t)$    & Accuracy at time $t$ \\
      $\text{Accuracy}_0$     & Accuracy at deployment \\
      $\lambda$               & Model sensitivity to drift \\
      $\mathcal{D}(\cdot\lVert\cdot)$ & Distribution divergence \\
    \end{tabular}
    }

    \vspace{0.15cm}
    \mlsysalert{Key:}{No code changed. The \emph{world} changed.}
  \end{column}
  \begin{column}{0.46\textwidth}
    \begin{mlsyscard}{errorstroke}
    \textbf{Example:}\\[0.1cm]
    {\footnotesize Recommendation system: 85\% accuracy at launch $\to$ 79.2\% after 6 months. No bugs --- just user behavior shift.\\[0.1cm]
    \textbf{Engineering response:} Set retraining triggers when $A(t)$ crosses a threshold.}
    \end{mlsyscard}
  \end{column}
\end{columns}

\end{frame}

% --- ACTIVE LEARNING: Quantitative Exercise ---
\begin{frame}{Exercise: When to Retrain?}
\note{
% [3 min]
% -- LINK: The Degradation Equation just introduced
% $\text{Accuracy}(t) \approx \text{Accuracy}_0 - \lambda \cdot \mathcal{D}$.
% This exercise applies it with concrete numbers.
%
% -- NARRATE: ``A fraud detector starts at 94\% accuracy. It drifts at
% $\alpha = 0.03$ per month with $\Delta$ growing linearly at 0.33 per month.
% The retraining trigger is 90\%. When does it fire?'' Give students 60 seconds.
%
% -- ENGAGE: Students calculate individually, then compare with a neighbor.
% Expected answer: $0.94 - 0.03 \times 0.33 \times t = 0.90$. Solving:
% $0.04 = 0.01 \times t$, so $t \approx 4$ months.
%
% -- WARN: Students may forget that $\Delta$ grows over time. Correct framing:
% drift is cumulative---each month the distribution moves further from training.
% IF STUCK: ``What is $A(t)$ at month 1? Month 2? See the pattern?''
%
% -- FLEX: [CORE] First quantitative exercise of the course.
% IF SHORT: Show the problem, reveal the answer after 30 seconds, move on.
}

\small
\begin{columns}[T]
  \begin{column}{0.55\textwidth}
    {\normalsize\bfseries When Does the Alarm Fire?}

    \vspace{0.2cm}
    A fraud detector has:
    \begin{itemize}\setlength\itemsep{0pt}
      \item $A_0 = 94\%$ accuracy at deployment
      \item Drift rate: $\alpha \cdot \Delta(t) \approx 1\%$ per month
      \item Retraining trigger: $A(t) < 90\%$
    \end{itemize}

    \vspace{0.15cm}
    \textbf{When does it cross the 90\% threshold?}

    {\footnotesize\textcolor{midgray}{(60 seconds --- then compare with a neighbor)}}
  \end{column}
  \begin{column}{0.42\textwidth}
    \pause
    \begin{mlsyscard}{errorstroke}
    \textbf{Solution:}\\[0.1cm]
    {\footnotesize
    $A(t) = 0.94 - 0.01t$\\[0.1cm]
    Month 1: 93\%\\
    Month 2: 92\%\\
    Month 3: 91\%\\
    \textcolor{errorstroke}{\textbf{Month 4: 90\%}} $\leftarrow$ trigger\\[0.1cm]
    Retrain by month 3--4 to stay above SLA.
    }
    \end{mlsyscard}
  \end{column}
\end{columns}

\end{frame}

% --- ACTIVE LEARNING: Predict ---
\begin{frame}{Predict: Which Constraint?}
\note{[2 min]
% -- LINK: Students just learned the Iron Law and the D·A·M taxonomy.
% This exercise applies that vocabulary before seeing the lighthouse models.
%
% -- NARRATE: ``A voice assistant must respond in under 200 ms. Which
% physical constraint dominates? Write your answer, then turn to a neighbor.''
%
% -- ENGAGE: Give 60 seconds, then cold-call 2--3 pairs.
% Expected answer: Network latency (if cloud-hosted, 50--100 ms per round
% trip), OR Power budget (if on-device, 100 mW budget for KWS model).
% Both are defensible---what matters is the reasoning, not the label.
%
% -- WARN: Many students will pick ``Compute Throughput'' because GPUs
% are marketed on TFLOPS. Correct framing: a 200 ms deadline on a
% 0.5 GFLOP/s microcontroller is compute-bound (KWS benchmark: 40 ms),
% but the binding system constraint is the power budget that forces the
% model onto a low-power device in the first place.
%
% -- FLEX: [CORE] First active learning moment. Prime for lighthouses.
% IF SHORT: Skip the neighbor discussion; cold-call one student immediately.
}

\centering
\vspace{0.8cm}
{\Large\bfseries Think--Write--Share}

\vspace{0.6cm}
{\large A voice assistant must respond in under 200 ms.\\[0.2cm]
\alert{What physical constraint dominates?}}

\vspace{0.4cm}
\begin{columns}[c]
  \begin{column}{0.22\textwidth}\centering\small Memory\\ Bandwidth\end{column}
  \begin{column}{0.22\textwidth}\centering\small Compute\\ Throughput\end{column}
  \begin{column}{0.22\textwidth}\centering\small Network\\ Latency\end{column}
  \begin{column}{0.22\textwidth}\centering\small Power\\ Budget\end{column}
\end{columns}

\vspace{0.5cm}
{\small\textcolor{lightgray}{Write your answer. Then turn to a neighbor.}}

\end{frame}

% --- Bridge: Frameworks to Lighthouses ---
\begin{frame}{From Diagnostic Tools to Patients}
\note{
% [1 min]
% -- LINK: Students just learned three analytical frameworks (D-A-M, Iron Law,
% Degradation Equation). This bridge slide connects them to the lighthouse
% models they will diagnose.
%
% -- NARRATE: ``Frameworks are diagnostic tools---like a stethoscope, a blood
% pressure cuff, and a thermometer. Lighthouses are the patients we diagnose.
% D-A-M tells you WHERE the bottleneck is. The Iron Law tells you HOW LONG
% it takes. The Degradation Equation tells you WHEN it will fail. Now we need
% patients to examine. That is what the five lighthouse models are.''
%
% -- ENGAGE: ``If D-A-M is the stethoscope, what is the Iron Law?'' Give 10
% seconds. Accept any reasonable medical analogy.
%
% -- WARN: Students may think frameworks and lighthouses are separate topics.
% Correct framing: they are inseparable---you always apply a framework TO a
% specific model on specific hardware.
%
% -- FLEX: [CORE] Essential bridge. Keep to 60 seconds.
% IF SHORT: State the analogy in one sentence and move on.
}

\centering
\vspace{0.6cm}
{\large\bfseries Frameworks are diagnostic tools.\\[0.3cm]
Lighthouses are the patients we diagnose.}

\vspace{0.5cm}
\begin{columns}[c]
  \begin{column}{0.30\textwidth}
    \centering
    \textcolor{datastroke}{\textbf{\DAM{}}}\\[0.1cm]
    {\small \emph{Where} is the\\bottleneck?}
  \end{column}
  \begin{column}{0.30\textwidth}
    \centering
    \textcolor{computestroke}{\textbf{Iron Law}}\\[0.1cm]
    {\small \emph{How long} does\\it take?}
  \end{column}
  \begin{column}{0.30\textwidth}
    \centering
    \textcolor{errorstroke}{\textbf{Degradation}}\\[0.1cm]
    {\small \emph{When} will\\it fail?}
  \end{column}
\end{columns}

\vspace{0.5cm}
{\normalsize $\downarrow$ \textbf{Apply these tools to five lighthouse models} $\downarrow$}

\end{frame}

% =============================================================================
\section{Lighthouse Models}
% =============================================================================

\begin{frame}{Five Lighthouse Models}
\note{[3 min]
% -- LINK: Students just saw the three analytical frameworks. The lighthouses
% are the five patients we will diagnose with those frameworks all semester.
%
% -- NARRATE: Walk through the diagram. ``Why do we need five, not one?
% Because different models stress different parts of the system. ResNet-50
% is compute-bound---it does a lot of arithmetic. GPT-2/Llama is
% memory-bandwidth-bound---it reads huge weight matrices every token.
% Each constraint regime requires a different engineering response.''
%
% -- ENGAGE: Ask: ``Which of these lighthouses runs on a smartphone?''
% Expected: MobileNetV2 (latency-bound). Follow-up: ``Which one would
% you find in a smart speaker?'' Expected: KWS (power-bound).
%
% -- WARN: Students will try to memorize which model maps to which
% bottleneck rather than understanding *why* the mapping follows from the
% architecture. The goal is not memorization---it is the reasoning pattern:
% architecture $\to$ arithmetic intensity $\to$ dominant Iron Law term.
%
% -- FLEX: [CORE] Essential for the whole course.
% IF SHORT: Name the five models and their bottleneck label only (30 seconds each).
}

% --- Layout: FULL-WIDTH IMAGE ---
\centering
\safeimg[width=0.92\textwidth,height=4.8cm,keepaspectratio]{lighthouse-models.pdf}

\vspace{0.1cm}
{\small Five models, five bottlenecks --- your \textbf{systems detectives} for the semester.}

\end{frame}

\begin{frame}{Lighthouse Models: Quick Reference}
\note{[2 min] Compact table for reference. Students will see these numbers
repeatedly. Emphasize: same Iron Law equation, different dominant term.}

\footnotesize
\renewcommand{\arraystretch}{1.2}
\begin{tabular}{@{}lp{2cm}lp{4cm}@{}}
  \toprule
  \textbf{Lighthouse} & \textbf{Bottleneck} & \textbf{Iron Law} & \textbf{Systems Question} \\
  \midrule
  \textbf{ResNet-50}    & \textcolor{computestroke}{Compute}  & $O/R_p$ & Is arithmetic fast enough? \\
  \textbf{GPT-2/Llama}  & \textcolor{datastroke}{Memory-BW}   & $D/BW$  & Can memory feed the accelerator? \\
  \textbf{MobileNetV2}  & \textcolor{routingstroke}{Latency}  & $L_{\text{lat}}$ & Can it meet real-time deadlines? \\
  \textbf{DLRM}         & \textcolor{datastroke}{Mem.\ Capacity} & $D/BW$  & Do terabyte-scale embedding tables fit in memory? \\
  \textbf{KWS}          & \textcolor{errorstroke}{Power}      & $P \times T$ & Can it run on a coin-cell battery? \\
  \bottomrule
\end{tabular}

\vspace{0.15cm}
{\footnotesize Same equation, different dominant term. Each lighthouse tests a different constraint regime.}

\end{frame}

% =============================================================================
\section{Course Roadmap}
% =============================================================================

\begin{frame}{Sixteen Chapters, Four Parts}
\note{[3 min]
% -- LINK: Students just saw the three frameworks and five lighthouses.
% This slide shows where each tool fits across the sixteen chapters.
%
% -- NARRATE: Walk through the roadmap. ``Foundations of ML Systems establishes
% the vocabulary and physics you will use all semester. Development and Training
% is where you build. Optimization and Acceleration is where you make it run
% efficiently. Deployment and Operations is where you ship it and keep it alive.
% Each part builds on the previous --- do not skip Part I thinking you know
% the basics. The D·A·M taxonomy and Iron Law from Part I are the lens for
% every chapter that follows.''
%
% -- ENGAGE: Ask: ``Which part do you think is hardest to engineer?''
% Give 15 seconds for a neighbor discussion. Expected responses: Part III
% (compression trade-offs), Part IV (ops and drift). Follow-up: ``Part I
% is often the most underestimated---the frameworks look simple but they
% recur in every single chapter at increasing depth.''
%
% -- FLEX: [CORE] Orientation. Keep to 3 minutes.
% IF SHORT: Walk the four parts only (30 seconds each), skip the engage question.
}

% --- Layout: FULL-WIDTH IMAGE ---
\centering
\safeimg[width=0.92\textwidth,height=4.8cm,keepaspectratio]{course-roadmap.pdf}

\vspace{0.1cm}
{\small \textbf{Foundations} $\to$ \textbf{Development \& Training} $\to$ \textbf{Optimization \& Acceleration} $\to$ \textbf{Deployment \& Operations} --- each part builds on the previous.}

\end{frame}

\begin{frame}{The Deployment Spectrum}
\note{[2 min] The gap between Cloud (H100, 80 GB) and TinyML (384 KB) is about
$10^5\times$ in memory and $10^6\times$ in compute---5 to 6 orders of magnitude,
not nine. You cannot ``shrink'' a cloud model to the edge --- each tier requires
fundamentally different engineering. This is the physical reality
that drives every design decision in the course.}

% --- Layout: FULL-WIDTH IMAGE ---
\centering
\safeimg[width=0.92\textwidth,height=4.8cm,keepaspectratio]{deployment-spectrum.pdf}

\vspace{0.1cm}
{\small $10^5\!\times$ memory, $10^6\!\times$ compute --- you cannot ``just shrink'' a cloud model.}

\end{frame}

\begin{frame}{Quick Check}
\note{[0.5 min] Micro-retrieval. Give students 30 seconds to recall.
This distributed practice improves retention (Roediger \& Karpicke, 2006).}

\centering
\Large\bfseries Quick check --- no peeking.\\[0.5cm]
\normalsize Which physical constraint prevents running GPT-4 on a smartphone?\\[0.3cm]
{\small 30 seconds --- then we continue.}
\end{frame}



% --- ACTIVE LEARNING 2: Discussion ---
\begin{frame}{Discussion: Where Does It Run?}
\note{[3 min] Turn-and-talk. Students discuss in pairs for 90 seconds.
This connects deployment spectrum to real decisions. Cold-call 2--3 pairs.
Multiple right answers exist --- the point is the reasoning, not the answer.}

\centering
\vspace{0.5cm}
{\large\bfseries Turn and Talk \textcolor{midgray}{(90 seconds)}}

\vspace{0.4cm}
{\large A hospital wants to deploy a model that detects\\
pneumonia from chest X-rays in real time.\\[0.2cm]
\alert{Where on the deployment spectrum should it run?}}

\vspace{0.35cm}
\begin{columns}[c]
  \begin{column}{0.22\textwidth}\centering\small\textcolor{computestroke}{\textbf{Cloud}}\\700 W\end{column}
  \begin{column}{0.22\textwidth}\centering\small\textcolor{datastroke}{\textbf{Edge}}\\15--40 W\end{column}
  \begin{column}{0.22\textwidth}\centering\small\textcolor{routingstroke}{\textbf{Mobile}}\\3--5 W\end{column}
  \begin{column}{0.22\textwidth}\centering\small\textcolor{errorstroke}{\textbf{TinyML}}\\100 mW\end{column}
\end{columns}

\vspace{0.2cm}
{\small\textcolor{lightgray}{What physical constraints drive your decision?}}

\end{frame}

% =============================================================================
\section{Skills \& Prerequisites}
% =============================================================================

\begin{frame}{What You'll Be Able To Do}
\note{[2 min] Concrete, measurable skills. These are the course-level learning
outcomes. Students should feel motivated --- these are real engineering skills.}

\footnotesize
By the end of this course, you will be able to:

\vspace{0.08cm}
\begin{enumerate}\setlength\itemsep{0pt}
  \item \textbf{Explain} why computational scale outperforms hand-coded expertise (the Bitter Lesson) and why ML systems fail silently
  \item \textbf{Apply} the \DAM{} Taxonomy to diagnose bottlenecks: data-bound, compute-bound, or machine-bound?
  \item \textbf{Calculate} memory budgets --- can this model fit on this device?
  \item \textbf{Design} serving pipelines with latency and throughput targets
  \item \textbf{Evaluate} hardware trade-offs across the deployment spectrum
  \item \textbf{Compress} models using quantization, pruning, and distillation
  \item \textbf{Monitor} production systems --- detect drift before users do
  \item \textbf{Make} principled design decisions grounded in physics, not hype
\end{enumerate}

\end{frame}

\begin{frame}{Prerequisites}
\note{[2 min] Be honest about what students need. This is a systems course ---
we assume OS and architecture basics. If students lack linear algebra,
recommend a refresher in week 1.}

\footnotesize
\begin{columns}[T]
  \begin{column}{0.48\textwidth}
    \textbf{Required background:}

    \vspace{0.1cm}
    \begin{itemize}\setlength\itemsep{0pt}
      \item \textbf{Operating systems} --- processes, memory, I/O
      \item \textbf{Computer architecture} --- caches, pipelines
      \item \textbf{Data structures} --- complexity analysis
      \item \textbf{Linear algebra} --- matrix multiply, norms
      \item \textbf{Probability} --- distributions, Bayes' rule
    \end{itemize}
  \end{column}
  \begin{column}{0.48\textwidth}
    \textbf{Helpful but not required:}

    \vspace{0.1cm}
    \begin{itemize}\setlength\itemsep{0pt}
      \item PyTorch or TensorFlow experience
      \item Systems programming (C/C++)
      \item Cloud computing basics
      \item Prior ML coursework
    \end{itemize}

    \vspace{0.05cm}
    \begin{mlsyscard}{datastroke}
    {\scriptsize \textbf{Note:} We teach ML concepts as needed. No ML expertise required to start.}
    \end{mlsyscard}
  \end{column}
\end{columns}

\end{frame}

\begin{frame}{Course Logistics}
\note{[3 min] Instructor should confirm grading breakdown, office hours,
collaboration policy, late policy, and resources for the current term.}

\small
\begin{columns}[T]
  \begin{column}{0.48\textwidth}
    \textbf{Grading} {\footnotesize(adjust per term)}

    \vspace{0.1cm}
    {\footnotesize
    \renewcommand{\arraystretch}{1.2}
    \begin{tabular}{@{}lr@{}}
      \toprule
      \textbf{Component} & \textbf{Weight} \\
      \midrule
      Problem Sets (5)        & 30\% \\
      Labs (4)                & 20\% \\
      Midterm Exam            & 20\% \\
      Final Project           & 25\% \\
      Participation           & 5\% \\
      \bottomrule
    \end{tabular}
    }
  \end{column}
  \begin{column}{0.48\textwidth}
    \textbf{Resources}

    \vspace{0.1cm}
    \begin{itemize}\setlength\itemsep{1pt}
      \item {\footnotesize \textbf{Textbook}: \textit{ML Systems}, Vol. I}
      \item {\footnotesize \textbf{Office hours}: See course site}
      \item {\footnotesize \textbf{Discussion forum}: See course site}
      \item {\footnotesize \textbf{Hardware labs}: Arduino, Raspberry Pi}
    \end{itemize}

    \vspace{0.15cm}
    \begin{mlsyscard}{routingstroke}
    {\footnotesize \textbf{Collaboration policy:} Discuss concepts freely. Write your own solutions. Cite all sources.}
    \end{mlsyscard}
  \end{column}
\end{columns}

\end{frame}

% =============================================================================
\section{A Taste}
% =============================================================================

\begin{frame}{Why Does GPT-4 Cost \$100M to Train?}
\note{[3 min] Compelling worked example. Walk through the back-of-envelope
calculation. The point: these are not magic numbers --- they follow from
the Iron Law. This is the kind of reasoning you will learn.
Ask: ``Where does the \$100M go?''}

\scriptsize
\textbf{A back-of-envelope calculation:}

\vspace{0.05cm}
\renewcommand{\arraystretch}{1.05}
\begin{tabular}{@{}llr@{}}
  \toprule
  \textbf{Factor} & \textbf{Estimate} & \textbf{Source} \\
  \midrule
  Parameters          & $\sim$1.8 trillion & Architecture choice \\
  Training FLOPs      & $\sim 2 \times 10^{25}$ & $\approx 6 \times N \times D_{\text{tokens}}$ \\
  Hardware            & $\sim$25,000 H100 GPUs & $R_{\text{peak}}$ per chip \\
  Training time       & $\sim$90 days & $O / (R_{\text{peak}} \cdot \eta \cdot N_{\text{GPUs}})$ \\
  Electricity         & $\sim$50 GWh & Power $\times$ time \\
  \midrule
  \textbf{Total cost} & \textbf{$\sim$\$100M} & Hardware + energy + ops \\
  \bottomrule
\end{tabular}

\vspace{0.05cm}
\begin{mlsyscard}{crimson}
{\scriptsize \textbf{Every number traces back to the Iron Law.} The cost follows from physics. This course teaches you to derive these numbers.}
\end{mlsyscard}

\end{frame}

\begin{frame}{What If You Had to Run It on a Phone?}
\note{[2 min] Dramatic contrast. GPT-4 needs 3.6 TB of memory at FP16
(1.8 trillion parameters $\times$ 2 bytes/param). A phone has 8 GB.
That is a 450$\times$ gap. Compression, quantization,
distillation --- these are the tools of Part III (Optimization and Acceleration).
Leave this as a teaser --- ``We will solve this problem.''}

\scriptsize
\begin{columns}[T]
  \begin{column}{0.55\textwidth}
    \textbf{GPT-4 at half precision (FP16):}
    Model size $\sim$3.6 TB. Phone RAM: 8 GB. \textcolor{errorstroke}{\textbf{Gap: 450$\times$.}}

    \vspace{0.05cm}
    \textbf{How do we close the gap?}
    \begin{itemize}\setlength\itemsep{0pt}
      \item \textcolor{computestroke}{\textbf{Quantization}} --- FP32 $\to$ INT4 (8$\times$ smaller)
      \item \textcolor{datastroke}{\textbf{Pruning}} --- remove redundant weights
      \item \textcolor{routingstroke}{\textbf{Distillation}} --- train a smaller model
    \end{itemize}
    Part III of this course teaches all three.
  \end{column}
  \begin{column}{0.42\textwidth}
    \renewcommand{\arraystretch}{1.0}
    \begin{tabular}{@{}lrr@{}}
      \toprule
      \textbf{Tier} & \textbf{RAM} & \textbf{Gap} \\
      \midrule
      Cloud (H100)   & 80 GB  & 45$\times$ \\
      Edge (Orin)    & 32 GB  & 112$\times$ \\
      Mobile         & 8 GB   & 450$\times$ \\
      TinyML         & 384 KB & 9.4M$\times$ \\
      \bottomrule
    \end{tabular}

    \vspace{0.05cm}
    \begin{mlsyscard}{errorstroke}
    You cannot ``just deploy'' a model. The deployment spectrum forces trade-offs at every tier.
    \end{mlsyscard}
  \end{column}
\end{columns}

\end{frame}


\begin{frame}{Muddiest Point}
\note{[1 min] One-minute paper. Collect responses (physical or digital).
Use responses to open the NEXT lecture with targeted clarification.
Angelo \& Cross (1993) --- powerful diagnostic tool.}

\centering
\Large\bfseries One-minute paper\\[0.5cm]
\normalsize Write down the \textbf{muddiest point} from today's lecture.\\[0.2cm]
{\small What concept was most confusing or unclear?}\\[0.5cm]
{\footnotesize\textcolor{midgray}{(Hand in on your way out --- or submit digitally)}}
\end{frame}

% --- Fallacies and Pitfalls ---
\begin{frame}{Fallacies and Pitfalls}
\note{
% [3 min]
% -- LINK: Students have now seen the five frameworks and five lighthouses.
% Before retrieval practice, surface the three most common Day 1 misconceptions
% so they leave knowing what NOT to believe.
%
% -- NARRATE: Read each fallacy aloud, then give the quantitative rebuttal.
% Pause 5 seconds between items to let the number land.
%
% -- ENGAGE: Ask after Fallacy 1: ``Has anyone here been burned by this?
% You optimized a model for weeks and it still failed to deploy?''
% Give 15 seconds for a show of hands.
%
% -- WARN: Students who came from ML research courses will push back hardest
% on Fallacy 2 (benchmark accuracy does not equal production readiness).
% That is healthy---let them push back and name the counter-evidence
% (94\% $\to$ 78--82\% is from the chapter; not invented).
%
% -- FLEX: [CORE] Required element---every chapter deck includes this.
% IF SHORT: Read Fallacies 1 and 2 only; skip Pitfall 3.
}

\footnotesize
\textbf{Fallacy:} \textit{Better algorithms automatically produce better systems.}

{\scriptsize A state-of-the-art Vision Transformer achieves 1--2\% higher accuracy than
ResNet-50 but requires 4$\times$ the FLOPs and 3$\times$ the memory bandwidth.
In production, a model that is 1\% more accurate but violates latency requirements
has effectively zero utility.}

\vspace{0.1cm}
\textbf{Fallacy:} \textit{High benchmark accuracy means the system is production-ready.}

{\scriptsize A sentiment model achieving 94\% accuracy on curated test data drops to
78--82\% in production as users employ slang, emojis, and context absent from benchmarks.
Distribution shift causes real-world degradation even when code is unchanged.}

\vspace{0.1cm}
\textbf{Pitfall:} \textit{Treating ML systems as traditional software that happens to include a model.}

{\scriptsize Traditional bugs produce stack traces within milliseconds; ML systems can silently
degrade 10--15\% over 3--6 months before anyone notices. CI/CD alone is not enough.}

\end{frame}

% --- ACTIVE LEARNING 3: Retrieval Practice ---
\begin{frame}{What Were the Key Ideas?}
\note{[2 min] Retrieval practice. Students write for 90 seconds, no notes.
Walk around the room. Do NOT show the next slide yet.
%
% -- FLEX: [CORE] First retrieval practice moment. Wait the full 90 seconds.
}

\centering
\vspace{1.5cm}
{\Large\bfseries Close your notes.}

\vspace{0.8cm}
{\large Write down the \textbf{4 most important ideas} from today.}

\vspace{0.8cm}
{\normalsize\textcolor{midgray}{90 seconds --- no peeking.}}

\end{frame}

\begin{frame}{Key Takeaways (1/2)}
\note{[1 min] Reveal first half. Walk through each bullet. Emphasize that every idea
introduced today will recur throughout the semester.
%
% -- ENGAGE: Ask students to compare what they wrote in the retrieval-practice
% slide with what they see here. How many of the items did they capture?
% Expected: 4--6 out of 11 is typical; fewer than 4 suggests the student was
% not tracking the through-line.
%
% -- WARN: If students recall fewer than 4 items, do NOT rush to show this
% slide---that defeats the retrieval practice. Give them the full 90 seconds
% first, then reveal as a check, not a rescue.
}

\small
\begin{itemize}\setlength\itemsep{1pt}
  \item \textbf{Software 2.0:} Weights are the new source code. Debugging moves from code to data.
  \item \textbf{Silent vs.\ loud failure:} ML systems degrade without crashing. Monitoring is not optional.
  \item \textbf{Bitter Lesson:} Scale beats hand-coded expertise. Systems that move compute efficiently determine AI progress.
  \item \textbf{Constraints drive architecture.} AI is infrastructure with physical laws --- not magic.
  \item \textbf{The 5\% problem:} Model code is 5\% of production. The other 95\% is systems engineering.
  \item \textbf{\DAM{} Taxonomy:} Data, Algorithm, Machine are interdependent. ``Which is the bottleneck?''
\end{itemize}

\end{frame}

\begin{frame}{Key Takeaways (2/2)}
\note{[1 min] Reveal second half. Connect all takeaways back to the course arc.
These frameworks and models recur in every chapter.}

\small
\begin{itemize}\setlength\itemsep{1pt}
  \item \textbf{Iron Law:} $T = D/BW + O/(R_p \cdot \eta) + L$. Slowest term dominates.
  \item \textbf{Degradation Eq.:} $\text{Accuracy}(t) \approx \text{Accuracy}_0 - \lambda \cdot \mathcal{D}(P_t\lVert P_0)$. Models fail silently as distributions drift.
  \item \textbf{Five Lighthouses:} ResNet-50, GPT-2/Llama, MobileNetV2, DLRM, KWS --- each isolates a bottleneck.
  \item \textbf{Deployment Spectrum:} Cloud to TinyML spans $10^5$--$10^6\!\times$ memory/compute. Different engineering per tier.
  \item \textbf{Course structure:} Foundations $\to$ Development \& Training $\to$ Optimization \& Acceleration $\to$ Deployment \& Operations.
\end{itemize}

\end{frame}

\begin{frame}{Next Lecture: Chapter 1 --- Introduction}
\note{[1 min] Forward hook. Chapter 1 goes deep on the D-A-M taxonomy,
the Iron Law with worked examples, and the five lighthouse models.
Assign pre-reading if applicable.}

\small
\centering
\vspace{0.5cm}
{\large\textbf{Chapter 1: Introduction}}

\vspace{0.3cm}
{\normalsize The Physics of AI Engineering}

\vspace{0.5cm}
\begin{columns}[T]
  \begin{column}{0.30\textwidth}
    \centering
    \textcolor{datastroke}{\textbf{\DAM{} Taxonomy}}\\[0.1cm]
    {\footnotesize Deep dive into all three axes}
  \end{column}
  \begin{column}{0.30\textwidth}
    \centering
    \textcolor{computestroke}{\textbf{Iron Law}}\\[0.1cm]
    {\footnotesize Worked examples with real hardware}
  \end{column}
  \begin{column}{0.30\textwidth}
    \centering
    \textcolor{errorstroke}{\textbf{Degradation}}\\[0.1cm]
    {\footnotesize Why models fail silently}
  \end{column}
\end{columns}

\vspace{0.5cm}
{\normalsize\textbf{Read:} Chapter 1 of the textbook before next class.}

\end{frame}


% =============================================================================
% BACKUP SLIDES
% =============================================================================
\appendix

\begin{frame}{Backup: Iron Law Quick Calculation}
\note{
% [3 min]
% -- NARRATE: Use if students want to see a concrete Iron Law application.
% Walk through the voice assistant example: data loading, compute, and
% orchestration overhead.
%
% -- FLEX: [OPTIONAL] Backup for students who want more practice.
}
\small
\textbf{Example: Voice assistant keyword detection on Arduino Nicla}

\vspace{0.15cm}
\renewcommand{\arraystretch}{1.15}
{\footnotesize
\begin{tabular}{@{}llll@{}}
  \toprule
  \textbf{Term} & \textbf{Variables} & \textbf{Calculation} & \textbf{Time} \\
  \midrule
  \textcolor{datastroke}{Data}    & 16 KB audio, 50 MB/s & $16\text{K}/50\text{M}$ & \textbf{0.32 ms} \\
  \textcolor{computestroke}{Compute} & 6M FLOPs, 0.5 GF, $\eta$=0.3 & $6\text{M}/(0.15\text{G})$ & \textbf{40 ms} \\
  \textcolor{errorstroke}{Overhead} & Wake-up + preprocess & --- & \textbf{2 ms} \\
  \midrule
  & & $T_{\text{total}}$ & \textbf{42.3 ms} \\
  \bottomrule
\end{tabular}
}

\vspace{0.15cm}
\textbf{Diagnosis:} Compute-bound. But power is the real constraint: 0.5 mW budget means duty cycling.
\end{frame}

\begin{frame}{Backup: Degradation Equation Month-by-Month}
\note{
% [3 min]
% -- NARRATE: Use if students struggled with the retraining exercise.
% Walk through the full table showing accuracy at each month.
%
% -- FLEX: [OPTIONAL] Scaffolding for the quantitative exercise.
}
\small
\textbf{Fraud detector degradation timeline:}

\vspace{0.15cm}
\renewcommand{\arraystretch}{1.15}
{\footnotesize
\begin{tabular}{@{}lrrrl@{}}
  \toprule
  \textbf{Month} & $\Delta(P)$ & $\alpha \cdot \Delta$ & $A(t)$ & \textbf{Status} \\
  \midrule
  0 & 0.00 & 0.0\% & 94.0\% & Deployed \\
  1 & 0.33 & 1.0\% & 93.0\% & Normal \\
  2 & 0.67 & 2.0\% & 92.0\% & Normal \\
  3 & 1.00 & 3.0\% & 91.0\% & \textcolor{routingstroke}{Watch} \\
  4 & 1.33 & 4.0\% & \textcolor{errorstroke}{90.0\%} & \textcolor{errorstroke}{\textbf{Retrain trigger}} \\
  6 & 2.00 & 6.0\% & 88.0\% & SLA violation \\
  \bottomrule
\end{tabular}
}

\vspace{0.15cm}
\textbf{Key:} Retrain at month 3--4. Waiting until month 6 means 2 months of SLA violations.
\end{frame}

\end{document}
